\documentclass[10pt,a4paper]{amsart}
\usepackage{amsmath,amssymb,amsthm,graphicx}
\usepackage{verbatim}
\newenvironment{reference}{\comment}{\endcomment}
%\newenvironment{reference}{}{}
\newenvironment{slogan}{\comment}{\endcomment}
\newenvironment{history}{\comment}{\endcomment}

% For commutative diagrams we use Xy-pic
\usepackage[all]{xy}

% We use 2cell for 2-commutative diagrams.
\xyoption{2cell}
% Theorem environments.
%
\theoremstyle{plain}
\newtheorem{theorem}[subsection]{Theorem}
\newtheorem{proposition}[subsection]{Proposition}
\newtheorem{lemma}[subsection]{Lemma}

\theoremstyle{definition}
\newtheorem{definition}[subsection]{Definition}
\newtheorem{example}[subsection]{Example}
\newtheorem{exercise}[subsection]{Exercise}
\newtheorem{situation}[subsection]{Situation}

\theoremstyle{remark}
\newtheorem{remark}[subsection]{Remark}
\newtheorem{remarks}[subsection]{Remarks}

\numberwithin{equation}{subsection}

% Macros
%
\def\lim{\mathop{\mathrm{lim}}\nolimits}
\def\colim{\mathop{\mathrm{colim}}\nolimits}
\def\Spec{\mathop{\mathrm{Spec}}}
\def\Hom{\mathop{\mathrm{Hom}}\nolimits}
\def\Ext{\mathop{\mathrm{Ext}}\nolimits}
\def\SheafHom{\mathop{\mathcal{H}\!\mathit{om}}\nolimits}
\def\SheafExt{\mathop{\mathcal{E}\!\mathit{xt}}\nolimits}
\def\Sch{\mathit{Sch}}
\def\Mor{\mathop{\mathrm{Mor}}\nolimits}
\def\Ob{\mathop{\mathrm{Ob}}\nolimits}
\def\Sh{\mathop{\mathit{Sh}}\nolimits}
\def\NL{\mathop{N\!L}\nolimits}
\def\CH{\mathop{\mathrm{CH}}\nolimits}
\def\proetale{{pro\text{-}\acute{e}tale}}
\def\etale{{\acute{e}tale}}
\def\QCoh{\mathit{QCoh}}
\def\Ker{\mathop{\mathrm{Ker}}}
\def\Im{\mathop{\mathrm{Im}}}
\def\Coker{\mathop{\mathrm{Coker}}}
\def\Coim{\mathop{\mathrm{Coim}}}

% Boxtimes
%
\DeclareMathSymbol{\boxtimes}{\mathbin}{AMSa}{"02}

%
% Macros for moduli stacks/spaces
%
\def\QCohstack{\mathcal{QC}\!\mathit{oh}}
\def\Cohstack{\mathcal{C}\!\mathit{oh}}
\def\Spacesstack{\mathcal{S}\!\mathit{paces}}
\def\Quotfunctor{\mathrm{Quot}}
\def\Hilbfunctor{\mathrm{Hilb}}
\def\Curvesstack{\mathcal{C}\!\mathit{urves}}
\def\Polarizedstack{\mathcal{P}\!\mathit{olarized}}
\def\Complexesstack{\mathcal{C}\!\mathit{omplexes}}
% \Pic is the operator that assigns to X its picard group, usage \Pic(X)
% \Picardstack_{X/B} denotes the Picard stack of X over B
% \Picardfunctor_{X/B} denotes the Picard functor of X over B
\def\Pic{\mathop{\mathrm{Pic}}\nolimits}
\def\Picardstack{\mathcal{P}\!\mathit{ic}}
\def\Picardfunctor{\mathrm{Pic}}
\def\Deformationcategory{\mathcal{D}\!\mathit{ef}}
\title{Cohomology of the spectrum of the base field}
\author{The Stacks Project Authors}
\date{stacks-project:weil.tex:local-snapshot-20261012; GNU FDL 1.2 or later}
\begin{document}
\maketitle
\noindent\textit{Editorial note.} All delivered mathematical text is the author's own, copied from the pinned Stacks Project chapter file weil.tex; no statement or proof step was rewritten, reordered or removed.\par\medskip
\noindent\textit{Editorial note.} The counted proof is the author's complete argument: the determination of the cohomology through the product formula with the base field, and the identification of the cycle class of the base with one through idempotence.\par\medskip
\noindent\textit{Editorial note.} The label lemma-base names the counted statement, and the proof environment immediately following it in the source is delivered verbatim as its proof.\par\medskip
\noindent\textit{Editorial note.} The counted statement is carried by the author's own lemma environment in the chapter source.\par\medskip
\noindent\textit{Editorial note.} Pre-open gate: python operations/agent-authoring/tools/precard\_gate.py --source weil.tex --label lemma-base --start 2770 --end 2797 --local-id 28155 returned RESULT=OK (L28155 free, occupancy FREE, public interval clear). Pre-assembly tag self-check: canonical tag resolved and not present in reference/public-items.jsonl.\par\medskip
\noindent\textit{Editorial note.} The statement and proof are copied from the pinned chapter file; the cited inputs are kept as external pointers with their canonical Stacks tags.\par\medskip
\noindent\textit{Allowed standard prerequisites (editorial).}\begin{itemize}
\item The cycle class of the spectrum of the base field is nonzero and the fundamental class is idempotent. Axiom (C)(c) and (C)(d) on compatibility with the intersection product and the behaviour of the class of the base field. Shows that the cycle class of the base is nonzero and idempotent, hence equal to one.
\item Behaviour of the theory under products with the base field. Axiom (B)(a) describing the cohomology of a product with the spectrum of the base field. Forces the cohomology of the base field to be one dimensional in degree zero.
\end{itemize}\medskip
\begin{lemma}
\label{lemma-base}
Assume given (D0), (D1), (D2), and (D3) satisfying (A), (B), and (C).
Then $H^i(\Spec(k)) = 0$ for $i \not = 0$ and there is a
unique $F$-algebra isomorphism $F = H^0(\Spec(k))$.
We have $\gamma([\Spec(k)]) = 1$ and $\int_{\Spec(k)} 1 = 1$.
\end{lemma}

\begin{proof}
By axiom (C)(d) we see that $H^0(\Spec(k))$ is nonzero and even
$\gamma([\Spec(k)])$ is nonzero.
Since $\Spec(k) \times \Spec(k) = \Spec(k)$ we get
$$
H^*(\Spec(k)) \otimes_F H^*(\Spec(k)) = H^*(\Spec(k))
$$
by axiom (B)(a) which implies (look at dimensions) that only
$H^0$ is nonzero and moreover has dimension $1$. Thus
$F = H^0(\Spec(k))$ via the unique $F$-algebra isomorphism
given by mapping $1 \in F$ to $1 \in H^0(\Spec(k))$.
Since $[\Spec(k)] \cdot [\Spec(k)] = [\Spec(k)]$ in the
Chow ring of $\Spec(k)$ we conclude that
$\gamma([\Spec(k)) \cup \gamma([\Spec(k)]) = \gamma([\Spec(k)])$
by axiom (C)(c). Since we already know that $\gamma([\Spec(k)])$ is nonzero
we conclude that it has to be equal to $1$.
Finally, we have $\int_{\Spec(k)} 1 = 1$ since
$\int_{\Spec(k)} \gamma([\Spec(k)]) = 1$
by axiom (C)(d).
\end{proof}

\end{document}
